%%% The main file. It contains definitions of basic parameters and includes all other parts.

%% Settings for single-side (simplex) printing
% Margins: left 40mm, right 25mm, top and bottom 25mm
% (but beware, LaTeX adds 1in implicitly)
\documentclass[12pt,a4paper]{report}
\setlength\textwidth{145mm}
\setlength\textheight{247mm}
\setlength\oddsidemargin{15mm}
\setlength\evensidemargin{15mm}
\setlength\topmargin{0mm}
\setlength\headsep{0mm}
\setlength\headheight{0mm}
% \openright makes the following text appear on a right-hand page
\let\openright=\clearpage

%% Settings for two-sided (duplex) printing
% \documentclass[12pt,a4paper,twoside,openright]{report}
% \setlength\textwidth{145mm}
% \setlength\textheight{247mm}
% \setlength\oddsidemargin{14.2mm}
% \setlength\evensidemargin{0mm}
% \setlength\topmargin{0mm}
% \setlength\headsep{0mm}
% \setlength\headheight{0mm}
% \let\openright=\cleardoublepage

%% Generate PDF/A-2u
\usepackage[a-2u]{pdfx}

%% Character encoding: usually latin2, cp1250 or utf8:
\usepackage[utf8]{inputenc}

%% Prefer Latin Modern fonts
\usepackage{lmodern}

%% Further useful packages (included in most LaTeX distributions)
\usepackage{amsmath}        % extensions for typesetting of math
\usepackage{amsfonts}       % math fonts
\usepackage{amsthm}         % theorems, definitions, etc.
\usepackage{bbding}         % various symbols (squares, asterisks, scissors, ...)
\usepackage{bm}             % boldface symbols (\bm)
\usepackage{graphicx}       % embedding of pictures
\usepackage{fancyvrb}       % improved verbatim environment
\usepackage{natbib}         % citation style AUTHOR (YEAR), or AUTHOR [NUMBER]
\usepackage[nottoc]{tocbibind} % makes sure that bibliography and the lists
			    % of figures/tables are included in the table
			    % of contents
\usepackage{dcolumn}        % improved alignment of table columns
\usepackage{booktabs}       % improved horizontal lines in tables
\usepackage{paralist}       % improved enumerate and itemize
\usepackage{xcolor}         % typesetting in color

%%% Basic information on the thesis

% Thesis title in English (exactly as in the formal assignment)
\def\ThesisTitle{Automated Program Minimization With Preserving of Runtime Errors}

% Author of the thesis
\def\ThesisAuthor{Denis Leskovar}

% Year when the thesis is submitted
\def\YearSubmitted{2021}

% Name of the department or institute, where the work was officially assigned
% (according to the Organizational Structure of MFF UK in English,
% or a full name of a department outside MFF)
\def\Department{Department of Distributed and Dependable Systems}

% Is it a department (katedra), or an institute (ústav)?
\def\DeptType{Department}

% Thesis supervisor: name, surname and titles
\def\Supervisor{doc. RNDr. Pavel Parízek, Ph.D.}

% Supervisor's department (again according to Organizational structure of MFF)
\def\SupervisorsDepartment{Department of Distributed and Dependable Systems}

% Study programme and specialization
\def\StudyProgramme{Computer Science}
\def\StudyBranch{System Programming}

% An optional dedication: you can thank whomever you wish (your supervisor,
% consultant, a person who lent the software, etc.)
\def\Dedication{%
Dedication.
}

% Abstract (recommended length around 80-200 words; this is not a copy of your thesis assignment!)
\def\Abstract{%
Debugging of large programs is a difficult and time consuming task. Given a runtime error, the
developer must first reproduce it. He then has to find the cause of the error and create a bugfix.
 This process can be made significantly more efficient by reducing the amount of code the developer
 has to look into. This paper introduces three different methodologies of automatically reducing
 a given program $P$ into its minimal runnable subset $P'$. The automatically generated program $P'$
 also has to result in the same runtime error as $P$. The main focus of the reduction is on correctness
 when operating in a concrete application domain set by this study. \par
Implementations of introduced methodologies written using the LLVM compiler infrastructure are then
 compared and classified. Performance is measured based on the statement count of the newly generated
 program and the speed at which the minimal variant was generated. Moreover, the limits of the three
 different approaches are investigated with respect to the general application domain. The paper
 concludes with an overview of the most general and efficient methodology.
}

% 3 to 5 keywords (recommended), each enclosed in curly braces
\def\Keywords{%
{automated debugging}, {code analysis}, {syntax tree}, {statement reduction}, {clang}
}

%% The hyperref package for clickable links in PDF and also for storing
%% metadata to PDF (including the table of contents).
%% Most settings are pre-set by the pdfx package.
\hypersetup{unicode}
\hypersetup{breaklinks=true}

% Definitions of macros (see description inside)
%%% This file contains definitions of various useful macros and environments %%%
%%% Please add more macros here instead of cluttering other files with them. %%%

%%% Minor tweaks of style

% These macros employ a little dirty trick to convince LaTeX to typeset
% chapter headings sanely, without lots of empty space above them.
% Feel free to ignore.
\makeatletter
\def\@makechapterhead#1{
  {\parindent \z@ \raggedright \normalfont
   \Huge\bfseries \thechapter. #1
   \par\nobreak
   \vskip 20\p@
}}
\def\@makeschapterhead#1{
  {\parindent \z@ \raggedright \normalfont
   \Huge\bfseries #1
   \par\nobreak
   \vskip 20\p@
}}
\makeatother

% This macro defines a chapter, which is not numbered, but is included
% in the table of contents.
\def\chapwithtoc#1{
\chapter*{#1}
\addcontentsline{toc}{chapter}{#1}
}

% Draw black "slugs" whenever a line overflows, so that we can spot it easily.
\overfullrule=1mm

%%% Macros for definitions, theorems, claims, examples, ... (requires amsthm package)

\theoremstyle{plain}
\newtheorem{thm}{Theorem}
\newtheorem{lemma}[thm]{Lemma}
\newtheorem{claim}[thm]{Claim}

\theoremstyle{plain}
\newtheorem{defn}{Definition}

\theoremstyle{remark}
\newtheorem*{cor}{Corollary}
\newtheorem*{rem}{Remark}
\newtheorem*{example}{Example}

%%% An environment for proofs

\newenvironment{myproof}{
  \par\medskip\noindent
  \textit{Proof}.
}{
\newline
\rightline{$\qedsymbol$}
}

%%% An environment for typesetting of program code and input/output
%%% of programs. (Requires the fancyvrb package -- fancy verbatim.)

\DefineVerbatimEnvironment{code}{Verbatim}{fontsize=\small, frame=single}

%%% The field of all real and natural numbers
\newcommand{\R}{\mathbb{R}}
\newcommand{\N}{\mathbb{N}}

%%% Useful operators for statistics and probability
\DeclareMathOperator{\pr}{\textsf{P}}
\DeclareMathOperator{\E}{\textsf{E}\,}
\DeclareMathOperator{\var}{\textrm{var}}
\DeclareMathOperator{\sd}{\textrm{sd}}

%%% Transposition of a vector/matrix
\newcommand{\T}[1]{#1^\top}

%%% Various math goodies
\newcommand{\goto}{\rightarrow}
\newcommand{\gotop}{\stackrel{P}{\longrightarrow}}
\newcommand{\maon}[1]{o(n^{#1})}
\newcommand{\abs}[1]{\left|{#1}\right|}
\newcommand{\dint}{\int_0^\tau\!\!\int_0^\tau}
\newcommand{\isqr}[1]{\frac{1}{\sqrt{#1}}}

%%% Various table goodies
\newcommand{\pulrad}[1]{\raisebox{1.5ex}[0pt]{#1}}
\newcommand{\mc}[1]{\multicolumn{1}{c}{#1}}

%%% Coding packages for pseudocode
\usepackage{algorithm} 
\usepackage{algpseudocode} 
\usepackage{listings}
\lstset
{
	frame=tlrb,
	numbers=left,
	breaklines=true,
	breakatwhitespace=true,
}

%%% Todo notes
\usepackage{xargs}	% Use more than one optional parameter in a new commands
\usepackage[colorinlistoftodos,prependcaption,textsize=tiny]{todonotes}
\newcommandx{\unsure}[2][1=]{\todo[linecolor=red,backgroundcolor=red!25,bordercolor=red,#1]{#2}}
\newcommandx{\change}[2][1=]{\todo[linecolor=blue,backgroundcolor=blue!25,bordercolor=blue,#1]{#2}}
\newcommandx{\info}[2][1=]{\todo[linecolor=OliveGreen,backgroundcolor=OliveGreen!25,bordercolor=OliveGreen,#1]{#2}}
\newcommandx{\improvement}[2][1=]{\todo[linecolor=Plum,backgroundcolor=Plum!25,bordercolor=Plum,#1]{#2}}
\newcommandx{\thiswillnotshow}[2][1=]{\todo[disable,#1]{#2}}

% Title page and various mandatory informational pages
\begin{document}
%%% Title page of the thesis and other mandatory pages

%%% Title page of the thesis

\pagestyle{empty}
\hypersetup{pageanchor=false}
\begin{center}

\centerline{\mbox{\includegraphics[width=166mm]{../img/logo-en.pdf}}}

\vspace{-8mm}
\vfill

{\bf\Large BACHELOR THESIS}

\vfill

{\LARGE\ThesisAuthor}

\vspace{15mm}

{\LARGE\bfseries\ThesisTitle}

\vfill

\Department

\vfill

{
\centerline{\vbox{\halign{\hbox to 0.45\hsize{\hfil #}&\hskip 0.5em\parbox[t]{0.45\hsize}{\raggedright #}\cr
Supervisor of the bachelor thesis:&\Supervisor \cr
\noalign{\vspace{2mm}}
Study programme:&\StudyProgramme \cr
\noalign{\vspace{2mm}}
Study branch:&\StudyBranch \cr
}}}}

\vfill

% Zde doplňte rok
Prague \YearSubmitted

\end{center}

\newpage

%%% Here should be a bound sheet included -- a signed copy of the "bachelor
%%% thesis assignment". This assignment is NOT a part of the electronic
%%% version of the thesis. DO NOT SCAN.

%%% A page with a solemn declaration to the bachelor thesis

\openright
\hypersetup{pageanchor=true}
\pagestyle{plain}
\pagenumbering{roman}
\vglue 0pt plus 1fill

\noindent
I declare that I carried out this bachelor thesis independently, and only with the cited
sources, literature and other professional sources. It has not been used to obtain another
or the same degree.

\medskip\noindent
I understand that my work relates to the rights and obligations under the Act No.~121/2000 Sb.,
the Copyright Act, as amended, in particular the fact that the Charles
University has the right to conclude a license agreement on the use of this
work as a school work pursuant to Section 60 subsection 1 of the Copyright~Act.

\vspace{10mm}

\hbox{\hbox to 0.5\hsize{%
In \hbox to 6em{\dotfill} date \hbox to 6em{\dotfill}
\hss}\hbox to 0.5\hsize{\dotfill\quad}}
\smallskip
\hbox{\hbox to 0.5\hsize{}\hbox to 0.5\hsize{\hfil Author's signature\hfil}}

\vspace{20mm}
\newpage

%%% Dedication

\openright

\noindent
\Dedication

\newpage

%%% Mandatory information page of the thesis

\openright

\vbox to 0.5\vsize{
\setlength\parindent{0mm}
\setlength\parskip{5mm}

Title:
\ThesisTitle

Author:
\ThesisAuthor

\DeptType:
\Department

Supervisor:
\Supervisor, \SupervisorsDepartment

Abstract:
\Abstract

Keywords:
\Keywords

\vss}

\newpage

\openright
\pagestyle{plain}
\pagenumbering{arabic}
\setcounter{page}{1}


%%% A page with automatically generated table of contents of the bachelor thesis

\tableofcontents

%%% Each chapter is kept in a separate file
\chapter*{Introduction}
\addcontentsline{toc}{chapter}{Introduction}

\change[inline]{TODO: Rewrite the~introduction to include goals.}

Automation of~routine tasks tied with software development has resulted in a~tremendous
increase in the~productivity of~software engineers. 

However, the~task of~debugging a~program
remains mostly manual chore. 
This is due to the~difficulty of~reliably encountering logic-based
runtime errors in the~code, a~task that, to this day, requires 
the developer's attention and supervision.

Let program $\mathcal{P}$ contain a~runtime error $E$ 
that consistently occurs when $\mathcal{P}$ is run 
with arguments $A$.
Since the~error $E$ is present at runtime and not compile-time,
it can be assumed that syntax wise the~code is mostly correct.
Therefore, any syntax-based error can be ruled out.
This, in turn, leaves us with a~set of~logical errors $\mathcal{E_L}$.
Those include wrongly indexed arrays and calculations that lead to either
the incorrect result or an altered control flow of~the program. 
Let $E \in \mathcal{E_L}$. As the~generality of~errors in $\mathcal{E_L}$
appears too complicated to be solved for all programming languages at once, 
it is necessary to break the~problem down for each programming language. 
This article is concerned with the~logical errors of~C and C++.
Although C is not a~subset of~C++, the~logical errors made in C can be
approached similarly to those in C++.
Both languages share mostly comparable constructs.
Finding the~cause of~a logical error in a~concrete language requires knowing
these constructs, their behavior, and their general handling.

To make finding the~cause of~an error a~systematic approach, one might try removing
unnecessary statements in the~code, thus minimizing the~program.
Let $\mathcal{P'}$ be a~minimal variant of~$\mathcal{P}$ such 
that $\mathcal{P'}$ results in the~same error $E$ as $\mathcal{P}$
when run with the~same arguments $A$.
If done carefully and correctly, $\mathcal{P'}$ represents the~smallest 
subset of~$\mathcal{P}$ regarding code
size, while preserving the~cause of~the error in that subset.
Upon manual inspection, the~developer is required to make less of~an effort to find
the cause in $\mathcal{P'}$ as opposed to $\mathcal{P}$.

The minimization of~a program can be achieved in numerous ways.
In further sections, the~article describes and compares three different approaches.
The first is based on naive statement removal and its consequences during runtime.
The second removes major chunks of~the code while
periodically testing the~generated program's correctness.
The third deploys a~sequence of~code
altering techniques, namely slicing and delta debugging.


\chapter{Title of the first chapter}

An~example citation: \cite{Andel07}

\section{Title of the first subchapter of the first chapter}

\section{Title of the second subchapter of the first chapter}

\chapter{Title of the second chapter}

\section{Title of the first subchapter of the second chapter}

\section{Title of the second subchapter of the second chapter}


\chapter{Conclusion}

%% \addcontentsline{toc}{chapter}{Conclusion}

Source code minimization is a computationally demanding search problem. 
Generally, to achieve optimal results, i.e., the global minimum, one must 
generate and validate all possible results. 
Such a task results in exponentially many validations and is thus not feasible 
for any non-minor input. 
We attempt to avoid as many validations as possible, improving the running 
time while preserving the optimal result.

Approaches shown in this paper significantly differ in their time complexity. 
Analysis shows how a naive exponential algorithm can be sped up using 
heuristics. Those heuristics include iterative deepening and validating 
dependencies. 
A combination of search techniques and static analysis helps us formulate 
a more refined naive algorithm. 
Moreover, a rough approximation of the optimal result can be achieved 
in polynomial time by deploying a simple binary search technique. 
The approximation is denoted as a local minimum since it does not share 
the same optimality properties as the global minimum.

Relevant preprocessing techniques were presented, and their performance impact 
was measured. 
Running static and dynamic slicing reduced the source code's size 
significantly while preserving the desired runtime error. 
Furthermore, both slicers also conserved the cause of the error.
This paper shows that combining the mentioned preprocessing techniques with 
the presented algorithms yields good reduction results in average use cases. 
Such cases include unstructured, structured, and object-oriented source code. 
Unstructured programs saw the best results; we suppose that might be due to 
the nature of our implementation. 
The execution time is dramatically reduced compared to naive minimization. 
Introduced heuristics shaved off an immense number of validations in 
the average case. 
Though, the complexity for generating optimal results remains exponential in 
the worst case.

The premise of this project was to find a sophisticated way of minimizing 
a program while preserving the desired runtime error. 
We found out the slicing-based technique worked the best on all but trivial 
inputs. 
This result has confirmed and validated our beliefs held while formulating 
this technique.

\section{Future work}

Suggested techniques and algorithms have the potential to work well. 
They, however, require reliable and easy-to-use implementations. 
As mentioned in Section~\ref{chap:limitations}, our implementation of 
a minimization tool is not user-friendly due to several limitations. 
Upcoming enhancements focus on removing or relaxing these limitations. 
The main focus is on supporting a more comprehensive range of inputs. 
This goal can be achieved by implementing multi-file input support, reducing 
programs that interact with the user, and supporting multi-threaded 
applications.

Ideas that have been explained but not implemented are another focus of 
future attention. 
For example, the analysis explains how a static analyzer can be utilized to 
achieve better results. 
However, due to technical reasons, the implementation of this step could not 
be finished. 
There is room for more complicated heuristics, instrumentation, or pattern 
recognition for further improvements to speed and accuracy.

AutoPIE - the implementation of this project - can only be launched on 
Unix-based and Unix-like systems. 
We plan to extend the support to other platforms by creating a Docker image. 
This way, the implementation can be quickly shipped and executed. 
In order to improve the user interface, we plan on introducing an extension 
for Visual Studio Code, through which AutoPIE can be launched.

In the last months of working on this project, we came across CReduce - 
a tool for test case size reduction. 
Like our implementation, CReduce uses techniques such as Delta debugging to 
reduce the size of a program while preserving a wanted property. 
We might analyze CReduce in the future and perhaps contribute our findings 
to the project.


%%% Bibliography
%%% Bibliography (literature used as a source)
%%%
%%% We employ bibTeX to construct the bibliography. It processes
%%% citations in the text (e.g., the \cite{...} macro) and looks up
%%% relevant entries in the bibliography.bib file.
%%%
%%% The \bibliographystyle command selects, which style will be used
%%% for references from the text. The argument in curly brackets is
%%% the name of the corresponding style file (*.bst). Both styles
%%% mentioned in this template are included in LaTeX distributions.

% \bibliographystyle{plainnat}    %% Author (year)
\bibliographystyle{unsrt}     %% [number]

\renewcommand{\bibname}{Bibliography}

%%% Generate the bibliography. Beware that if you cited no works,
%%% the empty list will be omitted completely.

\bibliography{bibliography}

%%% If case you prefer to write the bibliography manually (without bibTeX),
%%% you can use the following. Please follow the ISO 690 standard and
%%% citation conventions of your field of research.

% \begin{thebibliography}{99}
%
% \bibitem{lamport94}
%   {\sc Lamport,} Leslie.
%   \emph{\LaTeX: A Document Preparation System}.
%   2nd edition.
%   Massachusetts: Addison Wesley, 1994.
%   ISBN 0-201-52983-1.
%
% \end{thebibliography}


%%% Figures used in the thesis (consider if this is needed)
\listoffigures

%%% Tables used in the thesis (consider if this is needed)
%%% In mathematical theses, it could be better to move the list of tables to the beginning of the thesis.
\listoftables

%%% Abbreviations used in the thesis, if any, including their explanation
%%% In mathematical theses, it could be better to move the list of abbreviations to the beginning of the thesis.
\chapwithtoc{List of Abbreviations}

%%% Attachments to the bachelor thesis, if any. Each attachment must be
%%% referred to at least once from the text of the thesis. Attachments
%%% are numbered.
%%%
%%% The printed version should preferably contain attachments, which can be
%%% read (additional tables and charts, supplementary text, examples of
%%% program output, etc.). The electronic version is more suited for attachments
%%% which will likely be used in an electronic form rather than read (program
%%% source code, data files, interactive charts, etc.). Electronic attachments
%%% should be uploaded to SIS and optionally also included in the thesis on a~CD/DVD.
%%% Allowed file formats are specified in provision of the rector no. 72/2017.
\appendix
\chapter{Attachments}

\section{First Attachment}

\openright
\end{document}
